% GENERATED FILE. Edit canonical lessons, then run npm run generate:books.
% canonical-source-hash: fnv1a64:dca6a3e2a1670d2c
% canonical-lessons: HI-C187-girjaghar, HI-C187-thana, HI-C187-dakiya, HI-C187-sinema, HI-C187-park

\chapter{A church, A police station, A postman, The cinema, A park}
\label{ch:hi-a-church-a-police-station-a-postman-the-cinema-a-park}
\hlchaptermodality{187}

\begin{chapteropening}
\textbf{By the end of this chapter:} \emph{I can say a church, a police station, a postman, the cinema and a park.}

Complete the last lesson of chapter 187: I can say a church, a police station, a postman, the cinema and a park.
\end{chapteropening}

\section[\texorpdfstring{\hi{गिरजाघर} (girjāghar)}{गिरजाघर (girjāghar)}]{\hi{गिरजाघर} (girjāghar) — a church}

\label{lesson:HI-C187-girjaghar}



\begin{quote}
Before the new one: say the Hindi for hope, then the Hindi for a memory.
\end{quote}

\subsection*{You'll want to know: \hi{गिरजाघर}}
\textbf{\hi{गिरजाघर}} — \emph{girjāghar} — \textquotedblleft{}a church\textquotedblright{}.

\begin{cousinweb}
\textbf{\hi{गिरजाघर}} (\emph{girjāghar}) is \emph{girjā}, \textquotedblleft{}church\textquotedblright{}, usually traced to Portuguese \emph{igreja}, with \emph{ghar}, \textquotedblleft{}house\textquotedblright{}, from \emph{ghara} (CDIAL 4428).
\end{cousinweb}

\begin{grammarlens}[title={What you've built}]
One more word for this chapter.
\end{grammarlens}

\subsection*{Your turn}
\begin{itemize}
\raggedright
  \item \emph{Say it:} \emph{girjāghar}
  \item \emph{Say it:} \emph{girjāghar}, once more
  \item \emph{Say it:} say \emph{girjāghar}
  \item \emph{Recall:} say the Hindi for hope, then the Hindi for a memory, then say \emph{girjāghar} again
\end{itemize}

\begin{tcolorbox}[breakable,colback=teal!4,colframe=teal!35!black,title={Before you move on}]
What does \hi{गिरजाघर} mean? (\textquotedblleft{}A church\textquotedblright{}.) Say it once more.
\end{tcolorbox}
\section[\texorpdfstring{\hi{थाना} (thānā)}{थाना (thānā)}]{\hi{थाना} (thānā) — a police station}

\label{lesson:HI-C187-thana}



\begin{quote}
Before the new one: say the Hindi for a memory, then the Hindi for a church.
\end{quote}

\subsection*{You'll want to know: \hi{थाना}}
\textbf{\hi{थाना}} — \emph{thānā} — \textquotedblleft{}a police station\textquotedblright{}.

\begin{cousinweb}
\textbf{\hi{थाना}} (\emph{thānā}) is a tadbhava, a word inherited from Sanskrit and reshaped by sound change. It comes from \emph{sthāna}, \textquotedblleft{}firm stance; place of abode\textquotedblright{} (CDIAL 13753). CDIAL glosses the Hindi word \textquotedblleft{}police station\textquotedblright{}.
\end{cousinweb}

\begin{grammarlens}[title={What you've built}]
One more word for this chapter.
\end{grammarlens}

\subsection*{Your turn}
\begin{itemize}
\raggedright
  \item \emph{Say it:} \emph{thānā}
  \item \emph{Say it:} \emph{thānā}, once more
  \item \emph{Say it:} say \emph{thānā}
  \item \emph{Recall:} say the Hindi for a memory, then the Hindi for a church, then say \emph{thānā} again
\end{itemize}

\begin{tcolorbox}[breakable,colback=teal!4,colframe=teal!35!black,title={Before you move on}]
What does \hi{थाना} mean? (\textquotedblleft{}A police station\textquotedblright{}.) Say it once more.
\end{tcolorbox}
\section[\texorpdfstring{\hi{डाकिया} (ḍākiyā)}{डाकिया (ḍākiyā)}]{\hi{डाकिया} (ḍākiyā) — a postman}

\label{lesson:HI-C187-dakiya}



\begin{quote}
Before the new one: say the Hindi for a church, then the Hindi for a police station.
\end{quote}

\subsection*{You'll want to know: \hi{डाकिया}}
\textbf{\hi{डाकिया}} — \emph{ḍākiyā} — \textquotedblleft{}a postman\textquotedblright{}.

\begin{cousinweb}
\textbf{\hi{डाकिया}} (\emph{ḍākiyā}) is \emph{ḍāk}, \textquotedblleft{}the post\textquotedblright{}, with the ending \emph{-iyā}, \textquotedblleft{}one who deals with it\textquotedblright{}. CDIAL lists \emph{ḍāk}, \textquotedblleft{}relay, post\textquotedblright{}, under a reconstructed \emph{*ḍakka} (CDIAL 5519).
\end{cousinweb}

\begin{grammarlens}[title={What you've built}]
One more word for this chapter.
\end{grammarlens}

\subsection*{Your turn}
\begin{itemize}
\raggedright
  \item \emph{Say it:} \emph{ḍākiyā}
  \item \emph{Say it:} \emph{ḍākiyā}, once more
  \item \emph{Say it:} say \emph{ḍākiyā}
  \item \emph{Recall:} say the Hindi for a church, then the Hindi for a police station, then say \emph{ḍākiyā} again
\end{itemize}

\begin{tcolorbox}[breakable,colback=teal!4,colframe=teal!35!black,title={Before you move on}]
What does \hi{डाकिया} mean? (\textquotedblleft{}A postman\textquotedblright{}.) Say it once more.
\end{tcolorbox}
\section[\texorpdfstring{\hi{सिनेमा} (sinemā)}{सिनेमा (sinemā)}]{\hi{सिनेमा} (sinemā) — the cinema}

\label{lesson:HI-C187-sinema}



\begin{quote}
Before the new one: say the Hindi for a police station, then the Hindi for a postman.
\end{quote}

\subsection*{You'll want to know: \hi{सिनेमा}}
\textbf{\hi{सिनेमा}} — \emph{sinemā} — \textquotedblleft{}the cinema\textquotedblright{}.

\begin{cousinweb}
\textbf{\hi{सिनेमा}} (\emph{sinemā}) is a loan from English \emph{cinema}.
\end{cousinweb}

\begin{grammarlens}[title={What you've built}]
One more word for this chapter.
\end{grammarlens}

\subsection*{Your turn}
\begin{itemize}
\raggedright
  \item \emph{Say it:} \emph{sinemā}
  \item \emph{Say it:} \emph{sinemā}, once more
  \item \emph{Say it:} say \emph{sinemā}
  \item \emph{Recall:} say the Hindi for a police station, then the Hindi for a postman, then say \emph{sinemā} again
\end{itemize}

\begin{tcolorbox}[breakable,colback=teal!4,colframe=teal!35!black,title={Before you move on}]
What does \hi{सिनेमा} mean? (\textquotedblleft{}The cinema\textquotedblright{}.) Say it once more.
\end{tcolorbox}
\section[\texorpdfstring{\hi{पार्क} (pārk)}{पार्क (pārk)}]{\hi{पार्क} (pārk) — a park}

\label{lesson:HI-C187-park}



\begin{quote}
Before the new one: say the Hindi for a postman, then the Hindi for the cinema.
\end{quote}

\subsection*{You'll want to know: \hi{पार्क}}
\textbf{\hi{पार्क}} — \emph{pārk} — \textquotedblleft{}a park\textquotedblright{}.

\begin{cousinweb}
\textbf{\hi{पार्क}} (\emph{pārk}) is a loan from English \emph{park}.
\end{cousinweb}

\begin{grammarlens}[title={What you've built}]
One more word for this chapter.
\end{grammarlens}

\subsection*{Your turn}
\begin{itemize}
\raggedright
  \item \emph{Say it:} \emph{pārk}
  \item \emph{Say it:} \emph{pārk}, once more
  \item \emph{Say it:} say \emph{pārk}
  \item \emph{Recall:} say the Hindi for a postman, then the Hindi for the cinema, then say \emph{pārk} again
\end{itemize}

\begin{tcolorbox}[breakable,colback=teal!4,colframe=teal!35!black,title={Before you move on}]
What does \hi{पार्क} mean? (\textquotedblleft{}A park\textquotedblright{}.) Say it once more.
\end{tcolorbox}
